
\exo

On lance un dé à $6$ faces deux fois de suite et on note, dans l'ordre, les deux chiffres obtenus.

Pour chacune des affirmations suivantes, préciser si elle est vraie ou fausse. Justifier.
\begin{enumerate}
\item L'univers associé à cette expérience aléatoire est constitué de $36$ issues.
\item L'événement \og on obtient au moins un $6$ \fg{} est un événement élémentaire.
\item L'événement \og on obtient deux $6$ \fg{} est un événement élémentaire.
\end{enumerate}

\exo

Le tableau ci-dessous donne la probabilité de sortie de chaque face d'un dé cubique :
\begin{center}
\begin{tabular}{
>{\raggedright\arraybackslash}m{1.9cm}
*{6}{>{\centering\arraybackslash}m{1.1cm}}}
\hline
Issue & $1$ & $2$ & $3$ & $4$ & $5$ & $6$ \\ 
\hline
Probabilité & $0,25$ & $0,10$ & $x$ & $0,3$ & $0,2$ & $2x$ \\ 
\hline 
\end{tabular} 
\end{center}
\begin{enumerate}
\item
\begin{enumerate}
\item Le dé est-il équilibré ? Pourquoi ?
\item Calculer la valeur de $x$.
\end{enumerate}
\item Calculer la probabilité des événements suivants :
\begin{enumerate}
\item $A$ : \og On obtient un résultat pair. \fg{} 
\item $B$ : \og On obtient une valeur supérieure ou égale à $2$. \fg{} 
\item $C$ : \og On n'obtient pas plus de $3$. \fg{} 
\end{enumerate}
\end{enumerate}

\exo

On considère un dé cubique dont les faces sont numérotées de $1$ à $6$. Dans chacun des cas suivants, déterminer l'univers de l'expérience aléatoire décrite.
\begin{enumerate}
\item On lance le dé et on note le résultat obtenu.
\item On lance deux fois de suite le dé et on note la somme des résultats obtenus.
\item On lance deux fois de suite le dé et on note l'écart entre le résultat le plus grand et le plus petit.
\item On lance deux fois de suite le dé et on note le maximum des résultats obtenus.
\end{enumerate}

\exo

On a écrit sur des cartes les lettres suivantes : \clavier{L}~\clavier{I}~\clavier{E}~\clavier{G}~\clavier{E}~\clavier{A}~\clavier{R}~\clavier{D}.

Après avoir retourné les cartes, on en choisit une au hasard.
\begin{enumerate}
\item Établir la loi de probabilité de cette expérience aléatoire.
\item Calculer la probabilité d'obtenir une voyelle.
\item Calculer la probabilité d'obtenir une lettre qui n'apparaît qu'une fois dans le mot \textsc{LIEGEARD}.
\end{enumerate}

\exo

\begin{minipage}[t]{12cm}
Le tableau ci-contre donne la répartition des élèves d'une classe de Seconde. On choisit au hasard un élève de la classe et on s'intéresse aux événements :
\begin{itemize}
\item $G$ : \og l'élève choisi est un garçon. \fg{}
\item $E$ : \og l'élève choisi est externe. \fg{}
\end{itemize}
Combien d'issues possèdent les événements suivants ?
\vspace{-6pt}
\setlength{\columnsep}{1cm}
\setlength{\columnseprule}{0pt}
\begin{multicols}{4}
\begin{enumerate}
\item $G\cap E$.
\item $G\cup E$.
\item $\Gbarre \cap \Ebarre$.
\item $G\cup \Gbarre$.
\end{enumerate}
\end{multicols}
\vspace{-6pt}
\end{minipage}
\hfill
\begin{minipage}[t]{5.5cm}
~

\vspace{-12pt}
\setlength{\tabcolsep}{0.1cm}
\renewcommand{\arraystretch}{1.2}
\begin{tabular}{
>{\raggedright\arraybackslash}m{1.8cm}
*{2}{>{\centering\arraybackslash}m{1.3cm}}}
& Filles & Garçons \\ 
\hline
Externes & $5$ & $7$ \\ 
\hline 
Internes & $11$ & $12$ \\ 
\hline 
\end{tabular} 
\end{minipage}

\exo

\begin{minipage}[t]{12cm}
Sur le diagramme ci-contre, $\Omega$ est l'univers d'une expérience aléatoire, $A$ et $B$ sont deux événements. Chaque nombre indique le nombre d'issues pour la zone du diagramme correspondante.
\begin{enumerate}
\item Combien d'issues l'univers $\Omega$ possède-t-il ?
\item Combien d'issues possèdent les événements suivants ?
\vspace{-6pt}
\setlength{\columnsep}{1cm}
\setlength{\columnseprule}{0pt}
\begin{multicols}{4}
\begin{enumerate}
\item $A$.
\item $B$.
\item $\Abarre$.
\item $\Bbarre$.
\item $A\cap B$.
\item $A\cup B$.
\item $\overline{A\cap B}$.
\item $\overline{A\cup B}$.
\end{enumerate}
\end{multicols}
\vspace{-6pt}
\end{enumerate}
\end{minipage}
\hfill
\begin{minipage}[t]{5cm}
~

\vspace{-12pt}
\def\xmin{0}
\def\xmax{5}
\def\ymin{0}
\def\ymax{4}
\def\xunite{0.9cm}
\def\yunite{0.9cm}
\psset{unit=\xunite, xunit=\xunite, yunit=\yunite, algebraic=true, dimen=middle, dotstyle=*, dotsize=3pt 0, linewidth=0.8pt, arrowsize=3pt 2, arrowinset=0.25, comma=true}
\begin{pspicture*}(\xmin,\ymin)(\xmax,\ymax)
%\psgrid[xunit=\xunite, yunit=\yunite, subgriddiv=0, gridlabels=0, gridcolor=lightgray](0,0)(\xmin,\ymin)(\xmax,\ymax)
%\psaxes[labelFontSize=\scriptstyle, xAxis=true, yAxis=true, Dx=1, Dy=1, ticksize=-2pt 0, subticks=1]{->}(0,0)(\xmin,\ymin)(\xmax,\ymax)[{\footnotesize $x$},135][{\footnotesize $y$},-45]
\pspolygon(0,0)(4,0)(4,3)(0,3)
\psellipse(2.5,1.5)(1,0.5)
\pscircle(1.5,1.5){1}
\psline(3.5,3)(3.5,3.5)
\psline(1.5,2.5)(1.5,3.5)
\psline(3.5,1.5)(4.5,1.5)
\begin{footnotesize}
\uput{3pt}[090](3.5,3.5){$\Omega$}
\uput{3pt}[090](1.5,3.5){$A$}
\uput{3pt}[000](4.5,1.5){$B$}
\rput(1,1.5){$24$}
\rput(2,1.5){$15$}
\rput(3,1.5){$30$}
\rput(3,0.5){$20$}
\end{footnotesize}
\end{pspicture*}
\end{minipage}

\exo

Parmi les élèves de Première d'un lycée, \SI[retain-explicit-plus]{80}{\percent} suivent la spécialité Mathématiques, \SI[retain-explicit-plus]{40}{\percent} suivent la spécialité SES, et \SI[retain-explicit-plus]{30}{\percent} suivent les spécialités Mathématiques et SES.

On choisit au hasard un élève de Première, et on considère les événements $M$ : \og l'élève choisi suit la spécialité Mathématiques \fg{} et $S$ : \og l'élève choisi suit la spécialité SES \fg{}.
\begin{enumerate}
\item Comment est noté l'événement : \og l'élève choisi suit les spécialités Mathématiques et SES \fg{} ?

Quelle est sa probabilité ?
\item Décrire par une phrase l'événement $M\cup S$.

Quelle est sa probabilité ?
\end{enumerate}

\exo

Dans cet exercice, $A$ et $B$ désignent deux événements associés à une même expérience aléatoire.
\medskip
\begin{enumerate}
\item On sait que :
\[p(A)=0,1\qpvq p(\Bbarre)=0,6\qpvq p(A\cup B)=0,35.\]
Calculer $p(B)$ et $p(A\cap B)$.
\item On sait que :
\[p(A)=0,2\qpvq p(B)=0,5\qpvq p(A\cap B)=0,1.\]
Calculer $p(\overline{A\cup B})$.
\item On sait que :
\[p(\Abarre)=0,6\qpvq p(\Bbarre)=0,7\qpvq p(A\cup B)=0,55.\]
Calculer $p(\overline{A\cap B})$.
\item On sait que :
\[p(\Abarre)+p(\Bbarre)=0,9\qpvq p(A\cup B)=0,75.\]
Calculer $p(A\cap B)$.
\item On sait que :
\[p(A\cup B)=p(\overline{A\cap B})=0,85\qpvq p(A)=2\times p(B).\]
Calculer $p(A)$ et $p(B)$.
\end{enumerate}

\exo

Un jeu de $32$ cartes possède quatre \og couleurs \fg{} (c\oe{}ur, carreau, trèfle, pique). On tire successivement et avec remise deux cartes de ce jeu et on note, dans l'ordre, les couleurs obtenues.
\begin{enumerate}
\item Dénombrer les issues de cette expérience aléatoire (on pourra imaginer un arbre).
\item Calculer la probabilité de $A$ : \og on obtient deux trèfles \fg{}.
\item Calculer la probabilité de $B$ : \og on obtient deux cartes de même couleur \fg{}.
\item Calculer la probabilité de $C$ : \og on obtient au moins un carreau \fg{}.
\item Calculer la probabilité de $D$ : \og on obtient au plus un pique \fg{}.
\end{enumerate}

\exo

Le digicode d'une porte d'entrée est composé de $4$ chiffres allant de $0$ à $9$. On saisit au hasard un code à $4$ chiffres sur ce digicode.
\begin{enumerate}
\item Combien d'issues possède cette expérience ?
\item On considère les événements suivants :
\begin{itemize}
\item $A$ : \og Le code saisi commence par $5$. \fg{}
\item $B$ : \og Le code saisi ne contient pas de $0$. \fg{}
\end{itemize}
Donner deux issues réalisant chacun des événements suivants :
\vspace{-6pt}
\setlength{\columnsep}{1cm}
\setlength{\columnseprule}{0pt}
\begin{multicols}{4}
\begin{enumerate}
\item $A$.
\item $B$.
\item $\Abarre$.
\item $\Bbarre$.
\item $A\cap B$.
\item $A\cup B$.
\item $\Abarre\cap\Bbarre$.
\item $\Abarre\cup\Bbarre$.
\end{enumerate}
\end{multicols}
\vspace{-6pt}
\end{enumerate}

\exo

Une boîte contient un jeton bleu, un jeton jaune et un jeton rouge.

On tire un premier jeton, on note sa couleur, on le remet dans le sac puis on en tire un deuxième. Chaque jeton bleu rapporte 2 points, chaque jeton jaune rapporte $1$ point et chaque jeton rouge fait perdre $2$ points.
\begin{enumerate}
\item En utilisant un tableau à double entrée, déterminer les gains possibles à ce jeu (positifs ou négatifs).
\item À l'aide de ce tableau :
\begin{enumerate}
\item Déterminer le nombre de façons d'obtenir un gain nul.
\item Déterminer le nombre de façons d'obtenir un gain égal à $3$.
\end{enumerate}
\end{enumerate}

\exo

Le code d'ouverture d'une porte d'immeuble est constitué d'une lettre (choisie parmi A et B) suivie de deux chiffres (choisis parmi $2$, $3$ et $4$). Un code possible est B33.
\begin{enumerate}
\item À l'aide d'un arbre, dénombrer tous les codes possibles.
\item Dénombrer les codes ne comportant pas de chiffre pair.
\item Dénombrer les codes se terminant par un chiffre pair.
\item Quelle est la probabilité qu'un code choisi au hasard fasse apparaître deux fois le même chiffre ?
\end{enumerate}

\exo

On lance deux dés (un rouge et un bleu) dont les faces sont numérotées de $1$ à $6$. On note $r$ et $b$ les numéros des faces obtenues respectivement sur le dé rouge et sur le dé bleu.

La règle du jeu est la suivante :
\begin{itemize}
\item si $r$ est pair, on gagne $b$ euros ;
\item si $r$ vaut $1$ ou $3$, on gagne $b+1$ euros ;
\item si $r$ vaut $5$, on gagne $b+2$ euros.
\end{itemize}
\begin{enumerate}
\item Construire un tableau permettant de connaître le gain (en euros) selon les valeurs de $r$ et $b$.
\item Déterminer le nombre façons d'obtenir $8$ euros, puis la probabilité d'obtenir $8$ euros.
\item Déterminer la probabilité d'obtenir $4$ euros.
\end{enumerate}

\exo

Un sac contient deux jetons noirs ($N_1$ et $N_2$) et un jeton blanc ($B$), indiscernables au toucher.

On tire au hasard, successivement et avec remise, deux jetons du sac.
\begin{enumerate}
\item À l'aide d'un arbre, déterminer le nombre d'issues de cette expérience aléatoire.
\item On considère les événements $A$ : \og le premier jeton est blanc \fg{} et $B$ : \og le deuxième jeton est noir \fg{} .
\begin{enumerate}
\item Écrire $A$ et $B$ sous forme d'ensembles.

Combien d'issues réalisent chacun des événements $A$ et $B$ ?
\item Les événements $A$ et $B$ sont-ils incompatibles ?
\item Écrire sous forme d'ensembles les événements : \quad $\Abarre\qpvq \Bbarre \qpvq A\cup B\qpvq\Abarre\cap B$.
\end{enumerate}
\end{enumerate}

\exo 

Un jeu consiste à lancer deux dés tétraédriques équilibrés, dont les faces sont numérotées de $1$ à $4$. On gagne si la somme des deux valeurs obtenues est égale à $4$.

Ce jeu est simulé par la fonction Python ci-dessous :
\begin{pythoncodenumligne}
from math import randint
def jeu():
    a = ............
    b = ............
    if ..................:
        return "Gagné"
    else:
        return "Perdu"
\end{pythoncodenumligne}

\begin{enumerate}
\item Compléter la fontion \pythoninline{jeu}.
\item
\begin{enumerate}
\item Alice a joué \np{50} parties et gagné $13$ fois. Bob  a joué \np{1000} parties et gagné $178$ fois.

Quelle est la fréquence de victoire pour chacun d'eux ?
\item À l'aide d'un tableau à double entrée, montrer que ce jeu possède $16$ issues équiprobables, et déterminer la probabilité de gagner.
\item Commenter les résultats obtenus aux questions précédentes.
\end{enumerate}
\end{enumerate}

\exo

Une agence de voyage a effectué une enquête de satisfaction parmi \np{2500} clients, selon leur destination.

Les résultats sont résumés ci-dessous :
\begin{center}
\setlength{\tabcolsep}{0.1cm}
\renewcommand{\arraystretch}{1.2}
\begin{tabular}{
>{\raggedright\arraybackslash}m{4.5cm}
*{3}{>{\centering\arraybackslash}m{3cm}}}

&\textbf{À l'étranger} & \textbf{En France} & \textbf{Total} \\ 
\hline
\textbf{Clients satisfaits} & \np{1209} & \np{779} & \np{1988}\\ 
\hline 
\textbf{Clients non satisfaits} & \np{341} & \np{171} & \np{512}\\ 
\hline 
\textbf{Total} & \np{1550} & \np{950} & \np{2500}\\ 
\hline 
\end{tabular} 
\end{center}
On choisit au hasard un client de cette agence.
\begin{enumerate}
\item Déterminer la probabilité que le client soit satisfait.
\item Déterminer la probabilité que le client ne soit pas satisfait et qu'il ait voyagé en France.
\item On choisit un client ayant voyagé à l'étranger. Quelle est la probabilité qu'il ne soit pas satisfait ?
\item On choisit un client satisfait. Quelle est la probabilité qu'il ait voyagé en France ?
\end{enumerate}

\exo

On dispose d'un cube de bois dont l'arête mesure \SI{3}{\cm}.

On le peint en rouge, puis on le découpe en petits cubes dont les arêtes mesurent \SI{1}{\cm} :
\begin{center}
\psset{xunit=0.3cm,yunit=0.3cm,algebraic=true,dimen=middle,dotstyle=*,dotsize=2pt 0,linewidth=0.8pt,arrowsize=3pt 2,arrowinset=0.25}
\newcommand \xmin{-7}
\newcommand \xmax{4}
\newcommand \ymin{-0.5}
\newcommand \ymax{9.5}
\begin{pspicture*}(\xmin,\ymin)(\xmax,\ymax)
%\psgrid[xunit=1.0cm,yunit=1.0cm,subgriddiv=0,gridlabels=0,gridcolor=lightgray](0,0)(\xmin,\ymin)(\xmax,\ymax)
%\psaxes[labelFontSize=\scriptstyle,xAxis=true,yAxis=true,Dx=1,Dy=1,ticksize=-2pt 0,subticks=2]{->}(0,0)(\xmin,\ymin)(\xmax,\ymax)[{\footnotesize $x$},135][{\footnotesize $y$},-45]
\psline(-6.,7.)(0.,6.)
\psline(0.,6.)(3.,8.)
\psline(3.,8.)(-3.,9.)
\psline(-3.,9.)(-6.,7.)
\psline(-6.,7.)(-6.,1.)
\psline(3.,8.)(3.,2.)
\psline(-4.,6.6666)(-1.,8.6666)
\psline(-2.,6.3333)(1.,8.3333)
\psline(-5.,7.666)(1.,6.666)
\psline(-4.,8.3333)(2.,7.3333)
\psline(2.,7.3333)(2.,1.3333)
\psline(1.,6.666)(1.,0.6666)
\psline(0.,2.)(3.,4.)
\psline(0.,4.)(3.,6.)
\psline(0.,4.)(-6.,5.)
\psline(-6.,3.)(0.,2.)
\psline(-2.,0.3333)(-2.,6.3333)
\psline(-4.,6.6666)(-4.,0.6666)
\psline(-6.,1.)(0.,0.)
\psline(0.,0.)(3.,2.)
\psline(0.,0.)(0.,6.)
\end{pspicture*}
\end{center}

\begin{enumerate}
\item Combien de petits cubes obtient-on ainsi ?
\item On place tous ces petits cubes dans un sac, puis on en choisit un au hasard.

En justifiant soigneusement les réponses, déterminer les probabilités des événements suivants :
\begin{itemize}
\item $A$ : \og le petit cube choisi n'a aucune face peinte \fg{} ;
\item $B$ : \og le petit cube choisi a exactement une face peinte \fg{} ;
\item $C$ : \og le petit cube choisi a exactement quatre faces peintes \fg{} ;
\item $D$ : \og le petit cube choisi a au moins deux faces peintes \fg{}.
\end{itemize}
\end{enumerate}

\exo

Une puce est placée sur une droite graduée au point d'abscisse $0$. Chaque seconde, elle se déplace aléatoirement d'une unité vers la droite ou vers la gauche.
\begin{center}
\def\xmin{-4.5}
\def\xmax{4.5}
\def\ymin{-0.3}
\def\ymax{0.5}
\def\xunite{1.5cm}
\def\yunite{1.5cm}
\psset{unit=\xunite, xunit=\xunite, yunit=\yunite, algebraic=true, dimen=middle, dotstyle=*, dotsize=3pt 0, linewidth=0.8pt, arrowsize=3pt 2, arrowinset=0.25, comma=true}
\begin{pspicture*}(\xmin,\ymin)(\xmax,\ymax)
\psgrid[xunit=\xunite, yunit=\yunite, subgriddiv=0, gridlabels=0, gridcolor=lightgray](0,0)(\xmin,\ymin)(\xmax,\ymax)
\psaxes[labelFontSize=\scriptstyle, xAxis=true, yAxis=false, Dx=1, Dy=1, ticksize=-2pt 0, subticks=1]{->}(0,0)(\xmin,\ymin)(\xmax,\ymax)
\psdot[dotsize=5pt](0,0.2)
\psarc[arrows=<-](0.5,-0.4){0.8}{55}{120}
\end{pspicture*}
\end{center}

\begin{enumerate}
\item La puce effectue trois sauts.
\begin{enumerate}
\item Montrer que la puce ne peut pas finir en $+2$.
\item Calculer la probabilité qu'elle finisse en $+1$.
\end{enumerate}
\item La puce effectue quatre sauts.

Calculer la probabilité qu'elle effectue au plus trois sauts vers la droite.
\end{enumerate}

\exo

On choisit au hasard un jeton dans un sac contenant $20$ jetons blancs, $10$ jetons noirs et $n$ jetons rouges. 

Déterminer la valeur minimale de l'entier naturel $n$, de sorte que la probabilité de choisir un jeton rouge dépasse $0,99$.

